% =====================================================================
% DSAN 6650 Final Project: Research Paper (4-6 pages)
%
% Skeleton only. Every TODO describes what the section should contain and
% points to the matching part of docs/proposal.md.
%
% Build:   cd paper && latexmk -pdf main.tex
% Clean:   latexmk -c
%
% Open decision: the course may require a specific template (e.g. NeurIPS,
% IEEE, ACM). This file uses the plain article class until that is known;
% the section structure carries over to any template.
% =====================================================================
\documentclass[11pt]{article}

\usepackage[margin=1in]{geometry}
\usepackage[T1]{fontenc}
\usepackage{lmodern}
\usepackage{amsmath,amssymb}
\usepackage{booktabs}
\usepackage{graphicx}
\usepackage{xcolor}
\usepackage[numbers]{natbib}
\usepackage[hidelinks]{hyperref}

\graphicspath{{figures/}}

% Visible TODO marker; remove before submission.
\newcommand{\todo}[1]{\textcolor{red}{\textbf{TODO:} #1}}

\title{Comparative Study of Shallow and Deep Reinforcement Learning\\
Algorithms in Stock Trading with Performance-Based Weighted Ensemble}
\author{\todo{Team member names}\\DSAN 6650, Georgetown University}
\date{\todo{Submission date}}

\begin{document}
\maketitle

% ---------------------------------------------------------------------
\begin{abstract}
\todo{150--200 words. Problem (shallow vs.\ deep RL in a common historical
trading simulation), methods (Q-learning, SARSA, DQN, performance-weighted
ensemble), data and split, headline results for RQ1--RQ4, and one-sentence
takeaway. Write last.}
\end{abstract}

% ---------------------------------------------------------------------
\section{Introduction}
\label{sec:intro}
\todo{Motivation: trading as a controlled testbed for comparing RL
algorithms, not as a goal in itself. State the four research questions
(proposal Section 3) and list the contributions: unified historical-market
benchmark, shallow vs.\ deep comparison with multiple seeds and
out-of-sample testing, terminal-loss distribution, weighted ensemble.}

% ---------------------------------------------------------------------
\section{Background and Related Work}
\label{sec:background}
\todo{Brief MDP and TD-learning background; Q-learning~\cite{watkins1992q},
SARSA~\cite{rummery1994online}, DQN~\cite{mnih2015human}; see
\cite{sutton2018reinforcement}. A few sentences on prior RL-for-trading
work and on policy ensembles.}

% ---------------------------------------------------------------------
\section{Methods}
\label{sec:methods}

\subsection{Trading Environment}
\todo{Data, one stock per episode, daily steps, no leverage, no short
selling, transaction costs, observation vs.\ execution timing
(proposal Sections 5.1--5.2).}

\subsection{State and Action Spaces}
\todo{State features and how they are discretized for the tabular methods
(fit on training data only). The 9 discrete trade-sizing actions
(proposal Sections 5.3--5.4); Table~\ref{tab:actions} can move to an
appendix if space is tight.}

\begin{table}[h]
\centering
\caption{Discrete trade-sizing actions (tentative).}
\label{tab:actions}
\begin{tabular}{cl@{\hskip 2em}cl}
\toprule
Action & Operation & Action & Operation \\
\midrule
0 & Hold & & \\
1 & Buy with 25\% of cash & 5 & Sell 25\% of holdings \\
2 & Buy with 50\% of cash & 6 & Sell 50\% of holdings \\
3 & Buy with 75\% of cash & 7 & Sell 75\% of holdings \\
4 & Buy with 100\% of cash & 8 & Sell 100\% of holdings \\
\bottomrule
\end{tabular}
\end{table}

\subsection{Reward}
% Reconstructed from the outline; confirm with the team.
The reward is the per-step log return of portfolio value $V_t$, with
transaction costs reflected in $V_t$:
\begin{equation}
r_t = \ln\!\left(\frac{V_t}{V_{t-1}}\right), \qquad V_t > 0.
\label{eq:reward}
\end{equation}
\todo{State the shared discount factor $\gamma$ and its relation to
terminal return (proposal Section 5.5).}

\subsection{Algorithms}
\todo{Q-learning and SARSA update rules, $\epsilon$-greedy schedule; DQN
architecture, experience replay, target network. Note the shared training
budget in environment steps.}

\subsection{Performance-Based Weighted Ensemble}
% Reconstructed from the outline; confirm with the team.
Weights are a softmax over each policy's mean validation terminal return
$\bar{R}_i$ with temperature $\tau$:
\begin{equation}
w_i = \frac{\exp(\bar{R}_i/\tau)}{\sum_{j=1}^{N}\exp(\bar{R}_j/\tau)}.
\label{eq:weights}
\end{equation}
At each step the ensemble takes the action with the highest weighted vote,
\begin{equation}
a_t^{\mathrm{ens}} = \arg\max_{a\in\mathcal{A}} \sum_{i=1}^{N} w_i\,
\mathbf{1}\!\left[\pi_i(s_t)=a\right].
\label{eq:vote}
\end{equation}
\todo{Each member acts on the ensemble's own account state; tie-breaking
rule; equal-weight control ($w_i = 1/N$).}

% ---------------------------------------------------------------------
\section{Experimental Setup}
\label{sec:setup}
\todo{Stocks and data source; chronological split (train 2015--2021,
validation 2022, test 2023--2025 as three 12-month windows; tentative);
number of seeds; hyperparameters and how they were chosen (validation
only); baselines incl.\ buy-and-hold (proposal Sections 6--7).}

\paragraph{Metrics.}
Terminal return $R_T = (V_T - V_0)/V_0$ and terminal loss
$L_T = \max(0, -R_T)$, grouped into four categories (profit or break-even;
loss in $(0, 20\%]$; $(20\%, 40\%]$; $> 40\%$). Terminal loss is measured at
the end of each test window and is \emph{not} maximum drawdown.
\todo{Learning efficiency, stability, and trading-behavior metrics
(proposal Section 8).}

% ---------------------------------------------------------------------
\section{Results}
\label{sec:results}

\subsection{RQ1: Algorithm Performance}
\todo{Mean / median terminal return by method and test window.}

\begin{table}[h]
\centering
\caption{Out-of-sample terminal return (\%), mean $\pm$ std over seeds.
\todo{fill in}}
\label{tab:returns}
\begin{tabular}{lccc}
\toprule
Method & 2023 & 2024 & 2025 \\
\midrule
Q-learning & -- & -- & -- \\
SARSA & -- & -- & -- \\
DQN & -- & -- & -- \\
Weighted ensemble & -- & -- & -- \\
Equal-weight ensemble & -- & -- & -- \\
Buy-and-hold & -- & -- & -- \\
\bottomrule
\end{tabular}
\end{table}

\subsection{RQ2: Terminal Loss Distribution}
\todo{Counts per loss category (report zeros honestly) plus the continuous
return distribution.}

\begin{table}[h]
\centering
\caption{Terminal-loss category counts across test episodes. \todo{fill in}}
\label{tab:loss}
\begin{tabular}{lcccc}
\toprule
Method & Cat.~1 & Cat.~2 & Cat.~3 & Cat.~4 \\
\midrule
Q-learning & -- & -- & -- & -- \\
SARSA & -- & -- & -- & -- \\
DQN & -- & -- & -- & -- \\
Weighted ensemble & -- & -- & -- & -- \\
Equal-weight ensemble & -- & -- & -- & -- \\
Buy-and-hold & -- & -- & -- & -- \\
\bottomrule
\end{tabular}
\end{table}

\subsection{RQ3: Learning Efficiency and Stability}
\todo{Learning curves (return vs.\ environment steps); variability across
seeds, stocks, and years.}

\begin{figure}[h]
\centering
\fbox{\parbox[c][1.5in][c]{0.8\linewidth}{\centering
\todo{Learning curves: return vs.\ environment steps, mean $\pm$ std over
seeds.}}}
\caption{Learning curves. \todo{replace placeholder with
\texttt{\textbackslash includegraphics}}}
\label{fig:learning}
\end{figure}

\subsection{RQ4: Weighted Ensemble}
\todo{Weighted vs.\ equal-weight ensemble vs.\ best single algorithm on
return, terminal loss, and stability.}

\subsection{Trading Behavior}
\todo{Action-frequency breakdown, number of trades, fees, example
portfolio-value curves.}

% ---------------------------------------------------------------------
\section{Discussion and Limitations}
\label{sec:discussion}
\todo{Interpretation; limitations such as few stocks, single-asset
episodes, simulated execution, seeds on the same market window not being
independent samples, no claim of convergence to the optimal policy.}

% ---------------------------------------------------------------------
\section{Conclusion}
\label{sec:conclusion}
\todo{Answer RQ1--RQ4 in one or two sentences each; future work
(e.g.\ DQN on discrete states, more assets).}

% ---------------------------------------------------------------------
\bibliographystyle{plainnat}
\bibliography{references}

\end{document}
